The central risk is an incomplete candidate model. A proof for supported inputs
must state its dependencies and applicable namespace or rule subset. Unknown
features need a capture fallback or an explicitly incomplete outcome. Runtime
limits and failures reduce achieved coverage and must remain in the report.

A native renderer is also an experimental dependency. Recalculation, volatile
formulas, file-location functions and rendering variability can change output
without a transformation defect. A volatile dependency trace explains a supported
path rather than every image difference. A white exterior canvas difference
requires geometry and containment checks before classification. No global pixel
tolerance is inferred from one rare anti-aliasing case.

The Observatory reduces the effort of investigating these outcomes. Portable
reports and a future recorded-event replay can demonstrate execution and make
findings inspectable. They do not establish authenticity solely by being detailed,
and replay cannot supply missing historical telemetry. Scientific claims rely on
source identities, implementation versions and the independent evaluation.

\section{Conclusion}
DeltaSieve combines dependency-aware region selection with native rendering
and inspectable evidence. The bounded experiments demonstrate concrete
findings and fallbacks. Corpus-wide recall and matched capture cost remain
open.
